\begin{notechapter}
\chapter{Abstract Manifolds, Tangent and Cotangent Bundles}
\label{ch:tangent-bundles}
\chaptersubtitle{%
  When there is no ambient space}

\begin{chapterintro}
The set of straight lines of the plane, the set of angles, the set of
Gaussian distributions, the set of functions that a network can compute:
none of these comes inside an affine space. Each can still be described,
piece by piece, by $d$ real coordinates, and two descriptions of the same
piece are related by a diffeomorphism. This is all that
\cref{ch:submanifolds} really used, and it is taken here as the
definition. The price is that a point can no longer be moved by a vector
and that velocities can no longer be computed by subtracting points. What
survives is the tangent space at each point, now a vector space attached
to that point alone, its dual, and the rule that says how components
change with the coordinates: one rule for vectors, the inverse rule for
covectors. You will learn
\begin{itemize}[nosep]
  \item charts, atlases and manifolds, with examples that have no ambient
        space: the lines of the plane, the circle $\R/\Z$;
  \item tangent vectors as classes of curves, and why the tangent spaces
        at two points are unrelated;
  \item the differential of a map, the differential $\dd f(p)$ of a
        function as a cotangent vector, push-forward and pull-back;
  \item how the components of vectors, of covectors and of the Hessian at
        a critical point change with the chart;
  \item the tangent and cotangent bundles, vector fields and covector
        fields;
  \item what this says about reparametrizing a model and about the
        symmetries of its parameters.
\end{itemize}
\prerequisites{\cref{ch:duality}, \cref{ch:chain-rule-autodiff},
\cref{ch:inverse-implicit} and \cref{ch:submanifolds}. From this chapter
on, \emph{smooth} means infinitely differentiable; all manifolds, maps
and curves are smooth.}
\end{chapterintro}

\section{Charts and atlases}\label{tangent-bundles:sec:charts}

Two parametrizations of a submanifold differ by a diffeomorphism between
open subsets of $\R^d$ (\cref{submanifolds:cor:transition}). Reading the
parametrizations backwards, from the set to $\R^d$, gives the notion that
needs no ambient space.

\begin{definition}[Chart, atlas]\label{tangent-bundles:def:atlas}
Let $M$ be a set. A \emph{chart} of dimension $d$ on $M$ is a pair
$(U, \varphi)$ where $U\subset M$ and $\varphi:U\to\R^d$ is injective
with an open image. Two charts $(U, \varphi)$ and $(U', \varphi')$ are
\emph{compatible} if $\varphi(U\cap U')$ and $\varphi'(U\cap U')$ are
open and the \emph{transition map}
\[
  \tau = \varphi'\circ\varphi^{-1} :
  \varphi(U\cap U')\to\varphi'(U\cap U')
\]
is a diffeomorphism. An \emph{atlas} is a set of pairwise compatible
charts whose domains cover $M$.
\end{definition}

The components $x_1,\dots,x_d$ of $\varphi$ are functions on $U$, the
\emph{coordinates} of the chart; one writes $\varphi = (x_1,\dots,x_d)$.
The transition map expresses the coordinates $x'$ of one chart in terms
of the coordinates $x$ of the other (\cref{tangent-bundles:fig:charts}).

\begin{figure}[ht]
\centering
\begin{tikzpicture}[>=Stealth, scale=1.0]
  \begin{scope}
    \clip (-0.6,0.1) ellipse [x radius=1.6, y radius=1.0];
    \fill[bookaccent!18] (0.9,0.1) ellipse [x radius=1.6, y radius=1.0];
  \end{scope}
  \draw[thick] (0.15,0.1) ellipse [x radius=3.1, y radius=1.5];
  \draw[bookdef, thick] (-0.6,0.1) ellipse [x radius=1.6, y radius=1.0];
  \draw[bookpitfall, thick] (0.9,0.1) ellipse [x radius=1.6, y radius=1.0];
  \node at (-2.5,1.35) {$M$};
  \node[bookdef] at (-1.6,0.1) {$U$};
  \node[bookpitfall] at (1.95,0.1) {$U'$};
  \node[circle, fill, inner sep=1.2pt] at (0.15,0.25) {};
  \node[right] at (0.15,0.25) {$p$};
  % images
  \begin{scope}[yshift=-4.2cm, xshift=-4.6cm]
    \draw[bookdef, thick] (0,0) rectangle (2.8,1.8);
    \fill[bookaccent!18] (1.5,0.02) rectangle (2.78,1.78);
    \node[circle, fill, inner sep=1.2pt] at (2.0,1.0) {};
    \node[below] at (2.0,1.0) {$x$};
    \node[bookdef, below] at (1.0,0) {$\varphi(U)\subset\R^d$};
  \end{scope}
  \begin{scope}[yshift=-4.2cm, xshift=1.8cm]
    \draw[bookpitfall, thick] (0,0) rectangle (2.8,1.8);
    \fill[bookaccent!18] (0.02,0.02) rectangle (1.3,1.78);
    \node[circle, fill, inner sep=1.2pt] at (0.7,0.8) {};
    \node[below] at (0.7,0.8) {$x'$};
    \node[bookpitfall, below] at (1.8,0) {$\varphi'(U')\subset\R^d$};
  \end{scope}
  \draw[->, thick, bookdef] (-1.5,-0.95) -- (-2.9,-2.3);
  \node[bookdef, left] at (-2.25,-1.6) {$\varphi$};
  \draw[->, thick, bookpitfall] (1.8,-0.95) -- (3.2,-2.3);
  \node[bookpitfall, right] at (2.55,-1.6) {$\varphi'$};
  \draw[->, thick] (-1.6,-3.3) -- (1.6,-3.3);
  \node[above] at (0,-3.3) {$\tau = \varphi'\circ\varphi^{-1}$};
\end{tikzpicture}
\caption{Two charts on a set $M$. On the overlap $U\cap U'$ (shaded) a
point $p$ has two sets of coordinates, $x = \varphi(p)$ and
$x' = \varphi'(p) = \tau(x)$.}
\label{tangent-bundles:fig:charts}
\end{figure}

An atlas decides which subsets are open: $W\subset M$ is \emph{open} if
$\varphi(W\cap U)$ is open in $\R^d$ for every chart $(U, \varphi)$ of
the atlas.

\begin{proposition}\label{tangent-bundles:prop:topology}
Unions and finite intersections of open sets are open. The domain of
every chart is open, and a subset $W$ of the domain of a chart
$(U, \varphi)$ is open if and only if $\varphi(W)$ is open in $\R^d$.
\end{proposition}

\begin{proof}
Since $\varphi$ is injective, it commutes with unions and intersections
of subsets of $U$, which gives the first claim. The domain $U'$ of a
chart is open because $\varphi(U'\cap U)$ is open for every chart, by
compatibility. Let $W\subset U$. If $W$ is open then $\varphi(W)$ is
open, by definition. Conversely, if $\varphi(W)$ is open, then for any
chart $(U', \varphi')$ the set
$\varphi'(W\cap U') = \tau\bigl(\varphi(W)\cap\varphi(U\cap U')\bigr)$
is the image of an open set by a diffeomorphism, hence open.
\end{proof}

So each chart is a homeomorphism from its domain onto its image, and
continuity, limits and compactness make sense on $M$.

\begin{definition}[Manifold]\label{tangent-bundles:def:manifold}
A \emph{manifold} of dimension $d$ is a set $M$ with an atlas of
dimension $d$ such that
\begin{enumerate}[nosep]
  \item any two distinct points have disjoint open neighbourhoods (the
        \emph{Hausdorff} property);
  \item countably many charts of the atlas cover $M$.
\end{enumerate}
A chart compatible with all the charts of the atlas may be added to it;
two atlases whose union is an atlas define the same manifold.
\end{definition}

The two conditions exclude pathologies. Without the first, a sequence
could have two limits: take two copies of $\R$ and glue them along
$\R\setminus\{0\}$; the result has an atlas with two charts and two
origins that cannot be separated. The second will be needed to integrate
(\cref{ch:change-of-variables}).

\begin{proposition}[Submanifolds are manifolds]
\label{tangent-bundles:prop:submanifold}
Let $M$ be a smooth submanifold of dimension $d$ of an affine space
$\Aff$. The inverses $\varphi = \psi^{-1}$ of its parametrizations
(\cref{submanifolds:thm:four}\,(d)) form an atlas which makes $M$ a
manifold of dimension $d$, whose open sets are the sets $M\cap O$ with
$O$ open in $\Aff$.
\end{proposition}

\begin{proof}
The charts are compatible by \cref{submanifolds:cor:transition}. If
$\psi:\Omega\to M\cap U$ is a parametrization and $O$ is open in $\Aff$,
then $\varphi(M\cap O\cap U) = \psi^{-1}(O)$ is open; so the sets
$M\cap O$ are open, and $M$ is Hausdorff because $\Aff$ is. Conversely,
if $W\subset M$ is open and $p\in W$, then $\varphi(W\cap U)$ is open,
and since $\varphi$ is the restriction of a continuous map defined on an
open set of $\Aff$ (proof of \cref{submanifolds:prop:charts}), $W$
contains $M\cap B$ for a ball $B$ centred at $p$; $W$ is the intersection
of $M$ with the union of these balls. Finally, each point of $M$ lies in
a ball with rational radius, centred at a point with rational coordinates
in some frame, contained in the open set $U$ of a parametrization;
countably many such balls cover $M$.
\end{proof}

\begin{example}[Affine spaces]\label{tangent-bundles:ex:affine}
A frame $(o; e_1,\dots,e_n)$ of an affine space $\Aff$ gives a single
chart defined on all of $\Aff$, the affine coordinates. Two frames give
compatible charts, and the transition map is affine,
$x' = P^{-1}x + c$ (\cref{duality:prop:vector-change}). An affine space
is thus a manifold with a privileged family of global charts whose
transition maps are affine. An open subset of $\Aff$ is a manifold with
the same charts, restricted.
\end{example}

\begin{example}[The lines of the plane]\label{tangent-bundles:ex:lines}
Let $\mathcal L$ be the set of straight lines of $\R^2$. A non-vertical
line is $y = ax + b$ for a unique $(a, b)$; a non-horizontal line is
$x = a'y + b'$ for a unique $(a', b')$. This gives two charts,
$\varphi_1 = (a, b)$ on the set $U_1$ of non-vertical lines and
$\varphi_2 = (a', b')$ on the set $U_2$ of non-horizontal lines, each
with image $\R^2$. A line in $U_1\cap U_2$ has $a\neq0$, and
$y = ax + b$ reads $x = y/a - b/a$:
\[
  \tau(a, b) = \Bigl(\frac1a,\ -\frac ba\Bigr) ,
\]
a diffeomorphism of $\{a\neq0\}$ onto $\{a'\neq0\}$, equal to its own
inverse. The two charts form an atlas, and $\mathcal L$ is a manifold of
dimension $2$ (the Hausdorff property is checked in
\cref{tangent-bundles:exr:3}). The slope and intercept of a regression
line are the coordinates of the first chart; no single chart covers all
lines.
\end{example}

\begin{example}[The circle $\R/\Z$]\label{tangent-bundles:ex:circle}
Let $\mathcal C = \R/\Z$ be the set of classes $[x] = x + \Z$: the set of
angles, counted in turns. Each class has exactly one representative in
$(0, 1)$, except $[0]$, and exactly one in $(-\frac12, \frac12)$, except
$[\frac12]$. This gives two charts, $\varphi_1[x] = x\in(0, 1)$ and
$\varphi_2[x] = x\in(-\frac12, \frac12)$. Their domains overlap on two
arcs, and
\[
  \tau(x) = x \ \text{ for } 0 < x < \tfrac12 ,
  \qquad
  \tau(x) = x - 1 \ \text{ for } \tfrac12 < x < 1 .
\]
The transition map is a translation on each piece. The same construction
on $\R^2/\Z^2$ gives the torus, with four charts.
\end{example}

\begin{pitfall}[The same set can carry incompatible charts]
On the set $\R$, the chart $\varphi(x) = x$ and the chart
$\varphi'(x) = x^3$ are each an atlas, but they are not compatible:
$\varphi\circ\varphi'^{-1}(u) = u^{1/3}$ is not differentiable at $0$. A
manifold is a set \emph{with} an atlas. In the first, the function
$f(x) = x$ is smooth and has no critical point. In the second,
$f\circ\varphi'^{-1}(u) = u^{1/3}$: the same function is not even
differentiable. A change of parameters that is not a diffeomorphism
changes the calculus.
\end{pitfall}

\begin{definition}[Smooth map]\label{tangent-bundles:def:smooth}
Let $M$ and $N$ be manifolds. A map $F:M\to N$ is \emph{smooth} if it is
continuous and if, for all charts $(U, \varphi)$ of $M$ and $(W, \chi)$
of $N$, the \emph{local representative}
\[
  \chi\circ F\circ\varphi^{-1} :
  \varphi\bigl(U\cap F^{-1}(W)\bigr)\to\R^{\dim N}
\]
is smooth. A \emph{diffeomorphism} is a smooth bijection with a smooth
inverse.
\end{definition}

It is enough to check this for the charts of one atlas of each manifold,
around each point, because transition maps are diffeomorphisms.
Compositions of smooth maps are smooth. A function $f:M\to\R$ is smooth
when each $f\circ\varphi^{-1}$ is, and a curve is a smooth map from an
open interval into $M$. For submanifolds, this agrees with
\cref{submanifolds:def:c1} by \cref{submanifolds:prop:charts}.

\section{Tangent vectors}\label{tangent-bundles:sec:tangent}

On a manifold the quotient $(c(t) - c(0))/t$ has no meaning: points
cannot be subtracted. But whether two curves through $p$ have the same
velocity can still be decided, in a chart.

\begin{definition}[Tangent vector]\label{tangent-bundles:def:tangent}
Let $p\in M$. Two curves $c$ and $\tilde c$ with $c(0) = \tilde c(0) = p$
are \emph{tangent at $p$} if, in a chart $(U, \varphi)$ around $p$,
\[
  (\varphi\circ c)'(0) = (\varphi\circ\tilde c)'(0) .
\]
This is an equivalence relation. A class is a \emph{tangent vector} at
$p$, written $c'(0)$, and the set of classes is the \emph{tangent space}
$T_pM$.
\end{definition}

\begin{theorem}[The tangent space]\label{tangent-bundles:thm:tangent}
\begin{enumerate}[nosep]
  \item Tangency at $p$ does not depend on the chart.
  \item For each chart $(U, \varphi)$ around $p$, the map
        \[
          \theta_\varphi:T_pM\to\R^d , \qquad
          c'(0)\mapsto(\varphi\circ c)'(0) ,
        \]
        is a bijection, and for two charts
        $\theta_{\varphi'} = D\tau\bigl(\varphi(p)\bigr)\circ\theta_\varphi$.
  \item There is a unique structure of vector space on $T_pM$ for which
        every $\theta_\varphi$ is linear. Its dimension is $d$.
\end{enumerate}
\end{theorem}

\begin{proof}
Near $t = 0$, $\varphi'\circ c = \tau\circ(\varphi\circ c)$, so by the
chain rule
\[
  (\varphi'\circ c)'(0)
  = D\tau\bigl(\varphi(p)\bigr)\,(\varphi\circ c)'(0) .
\]
Since $D\tau(\varphi(p))$ is invertible, two curves have the same
velocity in the chart $\varphi$ if and only if they have the same
velocity in the chart $\varphi'$: this is (1), and the formula of (2).
The map $\theta_\varphi$ is injective by definition of the classes, and
surjective because, for $a\in\R^d$, the curve
$c(t) = \varphi^{-1}(\varphi(p) + ta)$ has $(\varphi\circ c)'(0) = a$.
For (3), transport the operations of $\R^d$ by $\theta_\varphi$. Another
chart gives the same operations, because
$\theta_{\varphi'}\circ\theta_\varphi^{-1} = D\tau(\varphi(p))$ is
linear.
\end{proof}

\begin{intuition}
A tangent vector is what all curves that leave $p$ ``in the same
direction at the same speed'' have in common. A chart turns it into a
column of $d$ numbers; another chart turns it into another column, and
the two columns are related by the Jacobian matrix of the change of
coordinates (\cref{tangent-bundles:fig:tangent}).
\end{intuition}

\begin{figure}[ht]
\centering
\begin{tikzpicture}[>=Stealth, scale=1.0]
  \begin{scope}
    \draw[thick] (0,0) ellipse [x radius=2.6, y radius=1.4];
    \node at (-2.1,1.2) {$M$};
    \draw[bookaccent, thick] (-1.8,-0.6) .. controls (-1.2,-0.6) and
      (-0.6,-0.36) .. (0,0) .. controls (0.6,0.36) and (1.2,0.5) ..
      (1.8,0.5);
    \draw[bookpitfall, thick] (-1.3,-1.15) .. controls (-0.9,-0.9) and
      (-0.5,-0.3) .. (0,0) .. controls (0.5,0.3) and (0.75,0.75) ..
      (0.9,1.15);
    \node[circle, fill, inner sep=1.3pt] at (0,0) {};
    \node[below right] at (0,-0.05) {$p$};
    \node[bookaccent, below] at (1.6,0.45) {$c$};
    \node[bookpitfall, left] at (0.8,1.0) {$\tilde c$};
  \end{scope}
  \draw[->, thick] (2.9,0) -- (4.3,0);
  \node[above] at (3.6,0) {$\varphi$};
  \begin{scope}[xshift=7.0cm]
    \draw[thick] (-2.4,-1.3) rectangle (2.4,1.3);
    \node at (1.7,-1.0) {$\varphi(U)$};
    \draw[bookaccent, thick] (-1.9,-0.9) .. controls (-1.2,-0.8) and
      (-0.5,-0.39) .. (0,0) .. controls (0.5,0.39) and (1.2,0.55) ..
      (1.9,0.6);
    \draw[bookpitfall, thick] (-1.4,-1.15) .. controls (-0.9,-1.0) and
      (-0.4,-0.31) .. (0,0) .. controls (0.4,0.31) and (0.8,0.8) ..
      (1.1,1.15);
    \draw[->, very thick] (0,0) -- (1.0,0.78);
    \node[circle, fill, inner sep=1.3pt] at (0,0) {};
    \node[below right] at (0,-0.05) {$\varphi(p)$};
    \node[above left] at (0.6,0.5) {$a$};
  \end{scope}
\end{tikzpicture}
\caption{Two curves through $p$ that are tangent at $p$: in a chart, they
have the same velocity $a = \theta_\varphi(c'(0))\in\R^d$. The class of
all such curves is one tangent vector.}
\label{tangent-bundles:fig:tangent}
\end{figure}

\begin{definition}[Coordinate basis]\label{tangent-bundles:def:basis}
For a chart $\varphi = (x_1,\dots,x_d)$ around $p$, the vectors
$\partial_i|_p = \theta_\varphi^{-1}(e_i)$, also written
$\frac{\partial}{\partial x_i}\big|_p$, form the \emph{coordinate basis}
of $T_pM$. The vector $\partial_i|_p$ is the velocity of the curve along
which $x_i$ increases at unit speed and the other coordinates stay
fixed. The \emph{components} of $v\in T_pM$ in the chart are the entries
of the column $[v] = \theta_\varphi(v)$: $v = \sum_iv_i\,\partial_i|_p$.
\end{definition}

\begin{proposition}[Tangent spaces of submanifolds and of affine spaces]
\label{tangent-bundles:prop:consistent}
Let $M$ be a submanifold of an affine space $\Aff$ with direction $V$.
The map that sends the class of a curve $c$ in $M$ to its velocity
$c'(0)\in V$ is a linear isomorphism from the tangent space of
\cref{tangent-bundles:def:tangent} onto the subspace $T_pM\subset V$ of
\cref{submanifolds:def:tangent}. It sends $\partial_i|_p$ to
$\partial_i\psi(t)$, where $\psi = \varphi^{-1}$ and $t = \varphi(p)$. In
particular, if $M$ is an open subset of $\Aff$, every tangent space
$T_pM$ is identified with the same vector space $V$.
\end{proposition}

\begin{proof}
A curve in $M$ is $c = \psi\circ(\varphi\circ c)$ near $0$, so its
velocity in $V$ is $D\psi(t)\,\theta_\varphi(c'(0))$. It depends only on
the class of $c$, and the map is $D\psi(t)\circ\theta_\varphi$: linear,
injective, with image $\im D\psi(t) = T_pM$
(\cref{submanifolds:thm:tangent}).
\end{proof}

\begin{keyidea}
On an affine space all tangent spaces are one vector space $V$: a vector
at $p$ can be compared with a vector at $q$, and added to any point. On a
manifold, $T_pM$ and $T_qM$ are different vector spaces with no given
map between them, and nothing like $p + v$ exists. For a submanifold the
ambient $V$ still contains them all, and the tangent flat $p + T_pM$ can
be drawn; without an ambient space there is no tangent flat, only
$T_pM$.
\end{keyidea}

\section{The differential of a map}\label{tangent-bundles:sec:differential}

\begin{definition}[Differential]\label{tangent-bundles:def:differential}
Let $F:M\to N$ be smooth and $p\in M$. The \emph{differential} of $F$ at
$p$ is the map
\[
  DF(p):T_pM\to T_{F(p)}N , \qquad c'(0)\mapsto(F\circ c)'(0) .
\]
It is also called the \emph{tangent map}, or the \emph{push-forward} of
tangent vectors by $F$.
\end{definition}

\begin{proposition}\label{tangent-bundles:prop:differential}
The differential is well defined and linear. In charts $\varphi$ around
$p$ and $\chi$ around $F(p)$, its matrix in the coordinate bases is the
Jacobian matrix of the local representative:
\[
  \theta_\chi\circ DF(p)
  = D\bigl(\chi\circ F\circ\varphi^{-1}\bigr)\bigl(\varphi(p)\bigr)
    \circ\theta_\varphi .
\]
For smooth maps $F:M\to N$ and $G:N\to P$,
$D(G\circ F)(p) = DG\bigl(F(p)\bigr)\circ DF(p)$.
\end{proposition}

\begin{proof}
Near $0$, $\chi\circ F\circ c
= (\chi\circ F\circ\varphi^{-1})\circ(\varphi\circ c)$, and the chain
rule gives
$(\chi\circ F\circ c)'(0)
= D(\chi\circ F\circ\varphi^{-1})(\varphi(p))\,(\varphi\circ c)'(0)$. So
the class of $F\circ c$ depends only on the class of $c$, and the
displayed formula holds; $DF(p)$ is linear as a composition of linear
maps. The chain rule is immediate from the definition:
$(G\circ F\circ c)'(0) = DG(F(p))\,(F\circ c)'(0)$.
\end{proof}

For submanifolds this is the differential of
\cref{submanifolds:prop:differential}, under the identification of
\cref{tangent-bundles:prop:consistent}. The theorems of
\cref{ch:inverse-implicit} are local and carry over through charts.

\begin{theorem}[Inverse function theorem on manifolds]
\label{tangent-bundles:thm:inverse}
Let $F:M\to N$ be smooth. If $DF(p)$ is invertible, there are open sets
$U\ni p$ and $W\ni F(p)$ such that $F$ is a diffeomorphism from $U$ onto
$W$.
\end{theorem}

\begin{proof}
The local representative $\chi\circ F\circ\varphi^{-1}$ has an invertible
differential at $\varphi(p)$, so it is a diffeomorphism between open
subsets of $\R^d$, with a smooth inverse
(\cref{inverse-implicit:thm:inverse,inverse-implicit:rem:remarks}).
Compose with the charts, which are diffeomorphisms onto their images.
\end{proof}

\begin{example}[The circle is the circle]\label{tangent-bundles:ex:s1}
The map $F[x] = (\cos2\pi x, \sin2\pi x)$ from $\R/\Z$ to the unit circle
$S^1\subset\R^2$ is well defined and bijective. In the charts of
\cref{tangent-bundles:ex:circle} and with the velocity in $\R^2$, its
differential sends $\partial_x$ to
$2\pi(-\sin2\pi x, \cos2\pi x)\neq0$, so it is invertible at every point:
$F$ is a diffeomorphism near each point, hence a diffeomorphism. The
abstract circle and the embedded circle are the same manifold; the first
description has no ambient space and its transition maps are
translations.
\end{example}

\section{Cotangent vectors}\label{tangent-bundles:sec:cotangent}

\begin{definition}[Cotangent space, differential of a function]
\label{tangent-bundles:def:cotangent}
The \emph{cotangent space} at $p$ is the dual $T_p^*M = (T_pM)^*$. For a
smooth function $f:M\to\R$, the \emph{differential} of $f$ at $p$ is the
cotangent vector
\[
  \dd f(p)\in T_p^*M , \qquad \dd f(p)\,c'(0) = (f\circ c)'(0) .
\]
\end{definition}

This is the differential of \cref{tangent-bundles:def:differential} for
$N = \R$, whose tangent spaces are all $\R$. Applied to the coordinates
themselves, it gives the dual of the coordinate basis.

\begin{proposition}[Components of a differential]
\label{tangent-bundles:prop:components}
Let $\varphi = (x_1,\dots,x_d)$ be a chart around $p$. The covectors
$\dd x_1(p),\dots,\dd x_d(p)$ form the basis of $T_p^*M$ dual to
$\partial_1|_p,\dots,\partial_d|_p$, and for a smooth function $f$
\[
  \dd f(p) = \sum_i\frac{\partial f}{\partial x_i}(p)\,\dd x_i(p) ,
  \qquad
  \frac{\partial f}{\partial x_i}(p)
  = \dd f(p)\,\partial_i|_p
  = \partial_i\bigl(f\circ\varphi^{-1}\bigr)\bigl(\varphi(p)\bigr) .
\]
\end{proposition}

\begin{proof}
Along the curve $c(t) = \varphi^{-1}(\varphi(p) + te_j)$, whose class is
$\partial_j|_p$, the coordinate $x_i$ is $x_i(p) + t\delta_{ij}$, so
$\dd x_i(p)\,\partial_j|_p = \delta_{ij}$. Along the same curve,
$f\circ c(t) = (f\circ\varphi^{-1})(\varphi(p) + te_j)$, whose derivative
at $0$ is the partial derivative of $f\circ\varphi^{-1}$. The expansion
of $\dd f(p)$ in the dual basis follows from
\cref{duality:thm:dual-basis}.
\end{proof}

The row of partial derivatives of $f$ in a chart is thus the row of
components of the covector $\dd f(p)$, as in
\cref{differential:prop:components}.

\begin{theorem}[Change of chart]\label{tangent-bundles:thm:change}
Let $x$ and $x'$ be the coordinates of two charts around $p$, and let
$J = \bigl[\partial x'_i/\partial x_j\bigr]$ be the Jacobian matrix of
the transition map at $\varphi(p)$. For $v\in T_pM$ and
$\alpha\in T_p^*M$,
\[
  [v]' = J\,[v] , \qquad [\alpha]' = [\alpha]\,J^{-1} ,
\]
and the number $\alpha(v) = [\alpha][v] = [\alpha]'[v]'$ does not depend
on the chart. The bases are related by
\[
  \partial'_j = \sum_i\frac{\partial x_i}{\partial x'_j}\,\partial_i ,
  \qquad
  \dd x'_i = \sum_j\frac{\partial x'_i}{\partial x_j}\,\dd x_j .
\]
\end{theorem}

\begin{proof}
$[v]' = \theta_{\varphi'}(v) = D\tau(\varphi(p))\theta_\varphi(v) = J[v]$
by \cref{tangent-bundles:thm:tangent}. This is the rule
$[v]' = P^{-1}[v]$ of \cref{duality:prop:vector-change} for the change of
basis matrix $P = J^{-1}$ from the basis $(\partial_i)$ to the basis
$(\partial'_j)$, whose columns hold the components of the $\partial'_j$;
and $J^{-1}$ is the Jacobian matrix $[\partial x_i/\partial x'_j]$ of
$\tau^{-1}$. Then $[\alpha]' = [\alpha]P = [\alpha]J^{-1}$ and
$\dd x'_i = \sum_j(P^{-1})_{ij}\,\dd x_j$ by
\cref{duality:prop:covector-change}.
\end{proof}

\begin{table}[ht]
\centering
\small
\begin{tabular}{@{}lll@{}}
\toprule
Object at $p$ & In the chart $x$ & In the chart $x' = \tau(x)$ \\
\midrule
tangent vector $v$ & column $[v]$ & $J[v]$ \\
cotangent vector $\alpha$, e.g.\ $\dd f(p)$ & row $[\alpha]$
  & $[\alpha]J^{-1}$ \\
pairing $\alpha(v)$ & $[\alpha][v]$ & the same number \\
Hessian at a critical point & $H$ & $J^{-\top}HJ^{-1}$ \\
inner product on $T_pM$ (\cref{ch:riemannian-optimization}) & $G$
  & $J^{-\top}GJ^{-1}$ \\
\bottomrule
\end{tabular}
\caption{How components change with the chart; $J$ is the Jacobian matrix
of $x\mapsto x'$. Compare with \cref{duality:tab:tests}, where the change
was linear and $J = P^{-1}$ was constant.}
\label{tangent-bundles:tab:change}
\end{table}

\begin{example}[Polar coordinates]\label{tangent-bundles:ex:polar}
On the plane without a closed half-line, take the charts $(r, \theta)$
and $(x, y)$, with $x = r\cos\theta$, $y = r\sin\theta$
(\cref{differential:ex:polar}). The Jacobian matrix of
$(r, \theta)\mapsto(x, y)$ is
$J = \begin{bsmallmatrix}\cos\theta & -r\sin\theta\\
\sin\theta & r\cos\theta\end{bsmallmatrix}$. Its columns give the polar
coordinate basis in terms of the Cartesian one, its rows give $\dd x$ and
$\dd y$:
\begin{align*}
  \partial_r &= \cos\theta\,\partial_x + \sin\theta\,\partial_y , &
  \dd x &= \cos\theta\,\dd r - r\sin\theta\,\dd\theta , \\
  \partial_\theta &= -r\sin\theta\,\partial_x + r\cos\theta\,\partial_y , &
  \dd y &= \sin\theta\,\dd r + r\cos\theta\,\dd\theta .
\end{align*}
For $f = x$, the row of partial derivatives in the polar chart is
$[\cos\theta\ \ {-r\sin\theta}]$. The vector $\partial_x$ has the polar
components $J^{-1}e_1 = (\cos\theta, -\frac1r\sin\theta)$. Row and column
agree only where $r = 1$ or $\sin\theta = 0$: in the polar chart the
transposed row of partial derivatives is not the vector $\partial_x$,
although it is in the Cartesian chart.
\end{example}

\begin{pitfall}[There is no gradient yet]
A manifold has differentials, which are covectors. To turn $\dd f(p)$
into a tangent vector one needs an inner product on $T_pM$, and a
manifold does not come with one. Reading the row of partial derivatives
as a column is choosing, silently, the inner product for which the
coordinate basis of the current chart is orthonormal; another chart makes
another choice. \Cref{ch:riemannian-optimization} makes the choice
explicit.
\end{pitfall}

\begin{definition}[Pull-back of a covector]\label{tangent-bundles:def:pullback}
Let $F:M\to N$ be smooth. The \emph{pull-back} by $F$ at $p$ is the dual
of the differential,
\[
  \dual{DF(p)}:T_{F(p)}^*N\to T_p^*M , \qquad
  \beta\mapsto\beta\circ DF(p) .
\]
\end{definition}

For a smooth function $g$ on $N$, the chain rule reads
$\dd(g\circ F)(p) = \dd g(F(p))\circ DF(p)$: the differential of
$g\circ F$ is the pull-back of the differential of $g$. Tangent vectors
are pushed forward and covectors are pulled back; in charts these are the
products of a Jacobian matrix with a column on its right and with a row
on its left, the two modes of \cref{chain-rule-autodiff:tab:modes}.

The second differential has no such general rule, because the second
derivative of a composition involves the first derivative
(\cref{gradient-hessian:rem:nonlinear}). At a critical point the extra
term vanishes.

\begin{proposition}[The Hessian at a critical point]
\label{tangent-bundles:prop:hessian}
Let $f:M\to\R$ be smooth and let $p$ be a critical point:
$\dd f(p) = 0$. There is a unique symmetric bilinear form $H_pf$ on
$T_pM$ such that
\[
  H_pf(v, v) = (f\circ c)''(0)
  \qquad\text{for every curve $c$ with $c'(0) = v$.}
\]
In a chart, its matrix is the matrix of second partial derivatives of
$f\circ\varphi^{-1}$ at $\varphi(p)$. Its signature does not depend on
the chart.
\end{proposition}

\begin{proof}
Let $g = f\circ\varphi^{-1}$, $x = \varphi(p)$ and
$\gamma = \varphi\circ c$. Differentiating
$(g\circ\gamma)'(t) = Dg(\gamma(t))\gamma'(t)$ once more,
\[
  (f\circ c)''(0) = D^2g(x)\bigl(\gamma'(0), \gamma'(0)\bigr)
  + Dg(x)\,\gamma''(0) ,
\]
and $Dg(x) = 0$ because $p$ is critical. So $(f\circ c)''(0)$ depends
only on $\gamma'(0) = \theta_\varphi(v)$, and it is the value at $[v]$ of
the quadratic form with matrix $[D^2g(x)]$. The left-hand side does not
mention the chart, and a quadratic form determines its symmetric bilinear
form. The signature is defined by \cref{spectral-svd:thm:sylvester}.
\end{proof}

At a point that is not critical, the term $Dg(x)\gamma''(0)$ remains, and
$\gamma''(0)$ can be changed at will by changing the chart: the matrix of
second partial derivatives is then not the matrix of any bilinear form on
$T_pM$. Whether a critical point is a non-degenerate minimum, a saddle,
or degenerate, is a property of $f$; the \emph{eigenvalues} of the matrix
$H$ are not, since $J^{-\top}HJ^{-1}$ has other eigenvalues unless $J$
is orthogonal.

\section{Bundles and fields}\label{tangent-bundles:sec:bundles}

\begin{definition}[Tangent and cotangent bundles]
\label{tangent-bundles:def:bundle}
The \emph{tangent bundle} $TM$ is the set of pairs $(p, v)$ with
$p\in M$ and $v\in T_pM$, with the projection $\pi(p, v) = p$. The
\emph{cotangent bundle} $T^*M$ is the set of pairs $(p, \alpha)$ with
$\alpha\in T_p^*M$.
\end{definition}

\begin{theorem}\label{tangent-bundles:thm:bundle}
The tangent bundle of a manifold of dimension $d$ is a manifold of
dimension $2d$. Each chart $(U, \varphi)$ of $M$ gives the chart
\[
  T\varphi:\pi^{-1}(U)\to\varphi(U)\times\R^d , \qquad
  (p, v)\mapsto\bigl(\varphi(p), \theta_\varphi(v)\bigr) ,
\]
and the transition maps are
$(x, a)\mapsto\bigl(\tau(x), D\tau(x)a\bigr)$. The same holds for $T^*M$,
with the transition maps
$(x, a)\mapsto\bigl(\tau(x), a\circ D\tau(x)^{-1}\bigr)$.
\end{theorem}

\begin{proof}
$T\varphi$ is a bijection onto the open set $\varphi(U)\times\R^d$ of
$\R^{2d}$ by \cref{tangent-bundles:thm:tangent}, which also gives the
transition maps. They are smooth, because $\tau$ is, with inverses of the
same form, and the images of the overlaps
$\varphi(U\cap U')\times\R^d$ are open. Countably many of these charts
cover $TM$. The projection $\pi$ is continuous, since
$T\varphi(\pi^{-1}(W)\cap\pi^{-1}(U)) = \varphi(W\cap U)\times\R^d$. Two
points of $TM$ over distinct points of $M$ are separated by the preimages
of disjoint open sets of $M$; two points over the same point lie in one
set $\pi^{-1}(U)$, which is homeomorphic to an open set of $\R^{2d}$. For
$T^*M$, use the components of covectors and
\cref{tangent-bundles:thm:change}.
\end{proof}

The differential of a smooth map $F$ is then a smooth map
$DF:TM\to TN$, $(p, v)\mapsto(F(p), DF(p)v)$: in charts it is
$(x, a)\mapsto(\hat F(x), D\hat F(x)a)$, with $\hat F$ the local
representative. A program that computes a value together with a
directional derivative computes $DF$.

\begin{definition}[Vector fields and covector fields]
\label{tangent-bundles:def:fields}
A \emph{vector field} is a smooth map $X:M\to TM$ with $X(p)\in T_pM$
for every $p$. A \emph{covector field} is a smooth map
$\alpha:M\to T^*M$ with $\alpha(p)\in T_p^*M$. In a chart,
\[
  X = \sum_iX_i\,\partial_i , \qquad \alpha = \sum_i\alpha_i\,\dd x_i ,
\]
and smoothness means that the component functions $X_i$, $\alpha_i$ are
smooth.
\end{definition}

The differential $\dd f:p\mapsto\dd f(p)$ of a smooth function is a
covector field, with components $\partial f/\partial x_i$. A covector
field applied to a vector field gives a function,
$\alpha(X)(p) = \alpha(p)X(p)$; in particular $\dd f(X)$ is the
derivative of $f$ along $X$.

\begin{proposition}[Pull-back of a covector field]
\label{tangent-bundles:prop:pullback}
Let $F:M\to N$ be smooth and $\beta$ a covector field on $N$. Then
\[
  (F^*\beta)(p) = \beta\bigl(F(p)\bigr)\circ DF(p)
\]
is a covector field on $M$, and $F^*(\dd g) = \dd(g\circ F)$ for every
smooth function $g$ on $N$.
\end{proposition}

\begin{proof}
In charts, with $\hat F$ the local representative, the row of components
of $F^*\beta$ at $x$ is $[\beta(\hat F(x))]\,[D\hat F(x)]$, a smooth
function of $x$. The second claim is the chain rule.
\end{proof}

\begin{pitfall}[Vector fields cannot be pushed forward]
Covector fields pull back along any smooth map. Vector fields do not
push forward: for $F(t) = t^2$ on $\R$ and $X = \partial_t$, the point
$1 = F(1) = F(-1)$ would receive both $DF(1)\partial_t = 2\partial_t$ and
$DF(-1)\partial_t = -2\partial_t$, and the points $s < 0$ nothing. A
vector field is pushed forward only by a diffeomorphism. This is why the
differential of the loss, a covector field, travels backwards through
every layer of a network, whatever the layer.
\end{pitfall}

\begin{example}[Flat, but not affine]\label{tangent-bundles:ex:flat}
On the circle $\R/\Z$ the transition maps are translations, with
Jacobian matrix $1$. So the coordinate vector $\partial_x$ is the same in
both charts: it is a vector field defined on the whole circle, without
zero, and $(p, a\,\partial_x)\mapsto(p, a)$ identifies $T(\R/\Z)$ with
the product $(\R/\Z)\times\R$. All tangent spaces are identified with
$\R$, as on an affine line, and $[x] + v = [x + v]$ moves every point by
a displacement $v$. Yet $\R/\Z$ is not an affine space: the displacement
from $[x]$ to $[y]$ is defined only up to an integer
(\cref{affine-spaces:exr:7}), and no chart covers it. The same holds for
the torus $\R^d/\Z^d$ (\cref{tangent-bundles:fig:flat}). On the sphere
$S^2$, by contrast, every vector field vanishes somewhere, a theorem of
topology that we do not prove; its tangent bundle is not a product.
\end{example}

\begin{figure}[ht]
\centering
\begin{tikzpicture}[>=Stealth, scale=1.0]
  \begin{scope}
    \draw[thick] (-2.6,1.6) -- (2.6,1.6);
    \foreach \x/\l in {-1.6/{x - 1}, 0.4/{x}, 2.4/{x + 1}} {
      \node[circle, fill=bookaccent, inner sep=1.4pt] at (\x,1.6) {};
      \node[above] at (\x,1.65) {$\l$};
    }
    \foreach \x in {-2,0,2} \draw (\x,1.52) -- (\x,1.68);
    \node[left] at (-2.6,1.6) {$\R$};
    \draw[->, thick] (0,1.3) -- (0,0.75);
    \draw[thick] (0,-1.0) circle [radius=1.0];
    \node[circle, fill=bookaccent, inner sep=1.4pt] at (0.309,-0.049) {};
    \node at (0.2,-0.4) {$[x]$};
    \draw[bookdef, very thick] (0,-1.0) ++(10:1.15) arc (10:350:1.15);
    \draw[bookpitfall, very thick] (0,-1.0) ++(-170:1.3) arc (-170:170:1.3);
    \node[bookdef, right] at (1.2,-1.75) {$\varphi_1$};
    \node[bookpitfall, left] at (-1.35,-1.75) {$\varphi_2$};
    \node at (0,-2.8) {the circle $\R/\Z$ and its two charts};
  \end{scope}
  \begin{scope}[xshift=6.3cm, yshift=-2.1cm]
    \draw[thick] (0,0) rectangle (3,3);
    \draw[->, thick, bookdef] (1.3,0) -- (1.7,0);
    \draw[->, thick, bookdef] (1.3,3) -- (1.7,3);
    \draw[->>, thick, bookpitfall] (0,1.3) -- (0,1.75);
    \draw[->>, thick, bookpitfall] (3,1.3) -- (3,1.75);
    \foreach \x in {0.6,1.5,2.4} \foreach \y in {0.5,1.4,2.3}
      \draw[->, bookaccent, thick] (\x-0.2,\y-0.1) -- (\x+0.2,\y+0.1);
    \node at (1.5,-0.7) {the torus $\R^2/\Z^2$};
  \end{scope}
\end{tikzpicture}
\caption{Left: the points $x - 1$, $x$, $x + 1$ of the line are one point
of the circle; each chart misses one point. Right: the torus is a square
whose opposite sides are identified. A constant vector field on the
square is a vector field on the torus, because the transition maps are
translations.}
\label{tangent-bundles:fig:flat}
\end{figure}

\section{Parameters as coordinates}\label{tangent-bundles:sec:ml}

\begin{mlthread}[Machine learning: reparametrizing a model]
The Gaussian distributions on $\R$ form a manifold of dimension $2$. The
pairs $(\mu, \sigma)$, $(\mu, \ln\sigma)$ and $(\mu, \sigma^2)$ are three
global charts on it, related by diffeomorphisms; a loss $L$ is a function
on the manifold and $\dd L$ a covector field.
\begin{itemize}[nosep]
  \item \emph{One covector, three rows.} At the same distribution the
        three charts give three rows of partial derivatives, related by
        $[\dd L]' = [\dd L]J^{-1}$. For the passage to
        $s = \ln\sigma$: $\partial L/\partial s
        = \sigma\,\partial L/\partial\sigma$.
  \item \emph{Three gradient steps.} Gradient descent in a chart reads
        the row as a column. A step $-\eta[\dd L]'^\top$ in the chart
        $x'$ is, to first order in $\eta$ and seen from the chart $x$,
        the step
        \[
          -\eta\,J^{-1}J^{-\top}[\dd L]^\top
          = -\eta\,(J^\top J)^{-1}[\dd L]^\top .
        \]
        Changing the chart preconditions the step by $(J^\top J)^{-1}$:
        each chart silently chooses the inner product with Gram matrix
        $J^\top J$ (\cref{tangent-bundles:fig:steps}). Training in
        $\ln\sigma$ multiplies the step on $\sigma$ by $\sigma^2$.
  \item \emph{What does not depend on the chart.} The critical points of
        $L$, and at each of them the signature of the Hessian
        (\cref{tangent-bundles:prop:hessian}). The eigenvalues of the
        Hessian at a minimum, often used to call a minimum ``sharp'' or
        ``flat'', change with the chart.
  \item \emph{Maps that are not charts.} Writing a weight as $w = u^3$,
        or a matrix as $W = BA$, is not a change of chart: the map from
        the new parameters is not a diffeomorphism. It creates critical
        points, $u = 0$ or $B = 0$, $A = 0$, that the loss does not
        have.
\end{itemize}
\end{mlthread}

\begin{figure}[ht]
\centering
\begin{tikzpicture}[>=Stealth, scale=2.6]
  \draw[->] (-0.7,0.5) -- (1.8,0.5) node[right] {$\mu$};
  \draw[->] (-0.6,0.4) -- (-0.6,2.4) node[above] {$\sigma$};
  \foreach \y in {1,2} {
    \draw (-0.63,\y) -- (-0.57,\y);
    \node[left] at (-0.63,\y) {$\y$};
  }
  \foreach \x in {0,1} {
    \draw (\x,0.47) -- (\x,0.53);
    \node[below] at (\x,0.47) {$\x$};
  }
  \node[circle, fill=bookdef, inner sep=1.8pt] at (0,1) {};
  \node[bookdef, below] at (0,0.95) {minimum};
  \draw[->, very thick, bookaccent] (1,2) -- (0.7,1.7);
  \draw[->, very thick, bookpitfall] (1,2) -- (0.7,1.098);
  \draw[->, very thick, gray] (1,2) -- (0.7,1.981);
  \node[circle, fill, inner sep=1.4pt] at (1,2) {};
  \node[right] at (1.03,2.02) {$(1, 2)$};
  \node[gray, above] at (0.6,2.03) {chart $(\mu, \sigma^2)$};
  \node[bookaccent, left] at (0.67,1.66) {chart $(\mu, \sigma)$};
  \node[bookpitfall, right] at (0.84,1.35) {chart $(\mu, \ln\sigma)$};
\end{tikzpicture}
\caption{One gradient step with the same step size $\eta = 1.2$ from the
same Gaussian distribution $(\mu, \sigma) = (1, 2)$, for the same loss
(the average negative log-likelihood of data with mean $0$ and variance
$1$), computed in three charts and drawn in the chart $(\mu, \sigma)$.
The covector $\dd L$ is the same; the three steps differ
(\cref{tangent-bundles:exr:5}).}
\label{tangent-bundles:fig:steps}
\end{figure}

\begin{mlthread}[Machine learning: symmetries and quotients]
Different parameters often give the same model.
\begin{itemize}[nosep]
  \item \emph{Examples.} A line $\alpha x + \beta y + \gamma = 0$ is
        unchanged when $(\alpha, \beta, \gamma)$ is multiplied by
        $c\neq0$: the manifold $\mathcal L$ of
        \cref{tangent-bundles:ex:lines} is the set of classes of
        parameters, and the map from parameters to lines is a submersion
        (\cref{tangent-bundles:exr:3}). In a network with a positively
        homogeneous activation, multiplying the weights into a hidden
        unit by $c > 0$ and those out of it by $1/c$ changes nothing;
        neither does permuting hidden units, nor replacing $(B, A)$ by
        $(BG^{-1}, GA)$ in a product $BA$.
  \item \emph{The differential of an invariant loss vanishes along the
        symmetry.} If $L(g\cdot\theta) = L(\theta)$ for a family of
        transformations $g_c$ depending smoothly on a parameter $c$,
        the velocity $\xi$ of the curve $c\mapsto g_c\cdot\theta$, taken
        where $g_c$ is the identity, satisfies $\dd L(\theta)\xi = 0$.
        For $L(w_1, w_2) = \ell(w_2w_1)$ with two scalar weights and the
        scaling $(cw_1, w_2/c)$, $\xi = (w_1, -w_2)$ and
        \[
          w_1\,\frac{\partial L}{\partial w_1}
          = w_2\,\frac{\partial L}{\partial w_2} .
        \]
        Along the gradient flow $\dot w = -[\dd L]^\top$ this gives
        $\frac{\dd}{\dd t}(w_1^2 - w_2^2) = 0$: the difference of the
        squared weights is conserved, and a gradient step of size $\eta$
        changes it by a term of order $\eta^2$ only
        (\cref{tangent-bundles:exr:9}).
  \item \emph{The Hessian is degenerate along the symmetry.} At a
        minimum, $H(\xi, \cdot) = 0$ for these directions $\xi$: the
        minima of the loss are not isolated in parameter space.
  \item \emph{The quotient.} The set of models is the set of classes of
        parameters. A general theorem, which we do not prove, makes it
        a manifold when the group of symmetries acts freely (only the
        identity fixes a point) and properly. At parameters with extra
        symmetry, such as two identical hidden units, the quotient is
        singular.
\end{itemize}
\end{mlthread}

\section{Where the chapter leads}\label{tangent-bundles:sec:outlook}

A manifold has tangent vectors, covectors, differentials, and a Hessian
at critical points. It has no gradient, no steepest direction, no notion
of a small step, and no way to move a point along a tangent vector.
\Cref{ch:riemannian-optimization} adds an inner product on each tangent
space, which gives gradients and lengths, and retractions, which replace
$p + v$. \Cref{ch:natural-gradient} shows that a family of
probability distributions carries one inner product that no
reparametrization can change. Covector fields come back in Part~V as the
objects that are integrated along curves.

\paragraph{Sources.} Loomis and Sternberg, \emph{Advanced Calculus},
Chapter~9: atlases and transition functions (\S9.1), manifolds as sets
with an equivalence class of atlases and the Hausdorff property
(Definition~3.1), differentiable maps (Definition~3.2), and tangent
vectors as classes of curves, with the change of chart and the induced
map on tangent spaces (\S9.4, equations (4.3) and (4.4)); they also
remark that on an open subset of a vector space the identity chart
identifies every tangent space with that vector space. Hubbard and
Hubbard stay with submanifolds of $\R^n$. N.~Boumal, \emph{An
Introduction to Optimization on Smooth Manifolds}, Cambridge University
Press, 2023, treats embedded submanifolds first (Chapter~3:
Definition~3.10, tangent spaces by curves in Definition~3.14 and as a
kernel in Theorem~3.15) and general manifolds in Chapter~8: charts and
atlases (Definitions~8.1 to 8.3), the atlas topology (Definition~8.17),
manifolds (Definition~8.21), tangent vectors as classes of curves
(Definition~8.33, Theorem~8.34), differentials (Definition~8.37) and the
tangent bundle (Theorem~8.43); his Chapter~9 treats quotient manifolds.
The conservation of the difference of squared norms between layers is
studied in S.~S.~Du, W.~Hu and J.~D.~Lee, Algorithmic regularization in
learning deep homogeneous models, NeurIPS, 2018; the dependence of
sharpness on the parametrization in L.~Dinh, R.~Pascanu, S.~Bengio and
Y.~Bengio, Sharp minima can generalize for deep nets, ICML, 2017.

\paragraph{Notebook.} The numerical claims of this chapter are checked
with PyTorch in \notebook{13-tangent-bundles.ipynb}.

\begin{chaptersummary}
\begin{itemize}[nosep]
  \item A manifold is a set with an atlas: charts into $\R^d$ whose
        transition maps are diffeomorphisms. Submanifolds, affine spaces,
        the lines of the plane and $\R/\Z$ are examples.
  \item A tangent vector at $p$ is a class of curves with the same
        velocity in a chart. $T_pM$ is a vector space of dimension $d$,
        attached to $p$; tangent spaces at different points are
        unrelated.
  \item The differential $DF(p)$ pushes tangent vectors forward; its
        matrix in coordinate bases is the Jacobian matrix of the local
        representative. The differential $\dd f(p)$ of a function is a
        cotangent vector, whose components are the partial derivatives.
  \item Under a change of chart with Jacobian matrix $J$: columns
        $[v]' = J[v]$, rows $[\alpha]' = [\alpha]J^{-1}$, and at a
        critical point $H' = J^{-\top}HJ^{-1}$.
  \item $TM$ and $T^*M$ are manifolds of dimension $2d$. Covector fields
        pull back along any smooth map; vector fields do not push
        forward.
  \item Gradient descent in a chart depends on the chart; critical points
        and the signature of the Hessian do not.
\end{itemize}
\end{chaptersummary}

\section*{Exercises}
\addcontentsline{toc}{section}{Exercises}

\begin{exercise}\label{tangent-bundles:exr:1}
On the unit circle $S^1\subset\R^2$ let
$\varphi_N(x, y) = x/(1 - y)$ on $S^1\setminus\{(0, 1)\}$ and
$\varphi_S(x, y) = x/(1 + y)$ on $S^1\setminus\{(0, -1)\}$. Show that
these two charts are compatible, with transition map $t\mapsto1/t$.
\solutionlink{tangent-bundles:exr:1}
\end{exercise}

\begin{exercise}\label{tangent-bundles:exr:2}
Let $M_1$ be $\R$ with the chart $x$, and $M_2$ be $\R$ with the chart
$x^3$. Show that the identity map $M_1\to M_2$ is smooth but not a
diffeomorphism, and that $x\mapsto x^{1/3}$ is a diffeomorphism from
$M_1$ onto $M_2$.
\solutionlink{tangent-bundles:exr:2}
\end{exercise}

\begin{exercise}\label{tangent-bundles:exr:3}
Let $\mathcal L$ be the manifold of lines of
\cref{tangent-bundles:ex:lines}.
\begin{enumerate}[nosep]
  \item Show that $\mathcal L$ has the Hausdorff property.
  \item Let $Q(\alpha, \beta, \gamma)$ be the line
        $\alpha x + \beta y + \gamma = 0$, for
        $(\alpha, \beta)\neq(0, 0)$. Compute the local representative of
        $Q$ in the chart $\varphi_1$, its Jacobian matrix, and the kernel
        of its differential.
\end{enumerate}
\solutionlink{tangent-bundles:exr:3}
\end{exercise}

\begin{exercise}\label{tangent-bundles:exr:4}
In the setting of \cref{tangent-bundles:ex:polar}, at the point with
$r = 2$ and $\theta = \pi/6$, compute the polar components of the vector
$v = \partial_y$ and of the covector $\alpha = \dd y$, and check that
$\alpha(v) = 1$ in both charts.
\solutionlink{tangent-bundles:exr:4}
\end{exercise}

\begin{exercise}\label{tangent-bundles:exr:5}
Let $L(\mu, \sigma) = \ln\sigma + \dfrac{1 + \mu^2}{2\sigma^2}$, the loss
of \cref{tangent-bundles:fig:steps}.
\begin{enumerate}[nosep]
  \item Compute the components of $\dd L$ at $(\mu, \sigma) = (1, 2)$ in
        the charts $(\mu, \sigma)$, $(\mu, s)$ with $s = \ln\sigma$, and
        $(\mu, v)$ with $v = \sigma^2$.
  \item Compute the point reached by one gradient step with $\eta = 1.2$
        in each chart, expressed in $(\mu, \sigma)$.
  \item Check the first-order formula of the text for the chart
        $(\mu, s)$.
\end{enumerate}
\solutionlink{tangent-bundles:exr:5}
\end{exercise}

\begin{exercise}\label{tangent-bundles:exr:6}
Let $F:\R/\Z\to\R/\Z$, $F[x] = [2x]$.
\begin{enumerate}[nosep]
  \item Show that $F$ is well defined and smooth, and that
        $DF(p)\,\partial_x = 2\,\partial_x$ at every point.
  \item Show that $F$ is a diffeomorphism near each point but not a
        diffeomorphism.
  \item Compute $F^*(\dd x)$.
\end{enumerate}
\solutionlink{tangent-bundles:exr:6}
\end{exercise}

\begin{exercise}\label{tangent-bundles:exr:7}
Let $F:\R\to\R^2$, $F(t) = (t^2, t^3)$, and let $\beta = x\,\dd y$ on
$\R^2$. Compute $F^*\beta$, and $DF(t)\,\partial_t$ in the coordinate
basis of $\R^2$. Where does $DF(t)$ fail to be injective?
\solutionlink{tangent-bundles:exr:7}
\end{exercise}

\begin{exercise}\label{tangent-bundles:exr:8}
On $\R^2$ consider the charts $(x, y)$ and
$(u, v) = (x + y + y^2, y)$.
\begin{enumerate}[nosep]
  \item For $f = x^2 - y^2$, which is critical at the origin, compute
        the matrices of second partial derivatives at the origin in both
        charts, check that $H' = J^{-\top}HJ^{-1}$, and compare the
        eigenvalues and the signatures.
  \item For $f = x$, compute the two matrices of second partial
        derivatives at the origin. Why is this not a contradiction?
\end{enumerate}
\solutionlink{tangent-bundles:exr:8}
\end{exercise}

\begin{exercise}\label{tangent-bundles:exr:9}
Let $L(w_1, w_2) = \frac12(w_2w_1 - 3)^2$.
\begin{enumerate}[nosep]
  \item Check that $w_1\,\partial L/\partial w_1
        = w_2\,\partial L/\partial w_2$ everywhere.
  \item Show that one gradient step of size $\eta$ changes
        $w_1^2 - w_2^2$ by
        $\eta^2\bigl((\partial L/\partial w_1)^2
        - (\partial L/\partial w_2)^2\bigr)$, and compute both sides for
        $(w_1, w_2) = (2, 0.5)$ and $\eta = 0.1$.
  \item Show that the Hessian of $L$ at a minimum is degenerate, and
        find its kernel.
\end{enumerate}
\solutionlink{tangent-bundles:exr:9}
\end{exercise}
\end{notechapter}
